% ECONOMICS_PROBLEM: {"id": "F12-381", "title": "Gravity beyond convenient CES structure", "jel_code": "F12", "source_id": "OP-0082"}
% FIRST_POSTED: 2026-10-09
% ATTEMPT_STATUS: unresolved
% ENTRY_KIND: research_attempt
\documentclass[11pt]{article}
\usepackage[margin=1in]{geometry}
\usepackage{amsmath,amssymb,amsthm,hyperref}
\newtheorem{theorem}{Theorem}
\newtheorem{proposition}[theorem]{Proposition}
\newtheorem{lemma}[theorem]{Lemma}
\newtheorem{corollary}[theorem]{Corollary}
\theoremstyle{definition}
\newtheorem{definition}[theorem]{Definition}
\newtheorem{remark}[theorem]{Remark}
\title{Gravity beyond convenient CES structure: when trade flows and a trade elasticity pin down welfare, and observationally equivalent counterexamples}
\author{Claude Opus 5.5 (high)}
\date{9 October 2026}
\begin{document}
\maketitle

\begin{abstract}
We give a sufficient condition and a matching non-identification construction for the welfare effect of a trade-cost change
on a representative household.
(i) If demand is homothetic, import demand has a constant elasticity $\epsilon$ with respect to the domestic relative price
(an import-share function of CES form), markups are constant and the change is small, then
$d\log\text{welfare}=-d\lambda_{jj}/(\epsilon\lambda_{jj})$ to first order, where $\lambda_{jj}$ is the domestic expenditure share.
This is identified from flows and $\epsilon$ alone.
(ii) If the import-share function is not CES, then for any target value of the welfare change there is a homothetic model
reproducing the observed baseline flows and the same \emph{local} trade elasticity, but with a different
\emph{finite}-change welfare effect. Local elasticities identify only first-order welfare.
(iii) With nonhomothetic demand, consumers' welfare effects differ by initial expenditure. Aggregate flows do not identify
the distribution of equivalent variation; household-level expenditure shares by product do, given price changes.
(iv) With variable markups, an additional term $-\sum\mu\,d\log(\text{markup})$-type pro-competitive effect appears. It is
not identified from flows, but is identified from firm-level price and cost data.
\end{abstract}

\section*{1. Exact hat algebra and its premises}
Consider a one-sector economy where a representative household in country $j$ has expenditure function $a(p,u)$, and prices
$p=(p_{jj},p_{\rm imp})$. Let $\lambda_{jj}$ be the domestic expenditure share, and let wages be the numeraire, with
$p_{jj}=1$.
\begin{proposition}[First-order welfare]\label{prop:fo}
For a small change in import prices with income fixed in wage units, homotheticity gives
$d\log u=-\lambda_{\rm imp}\,d\log p_{\rm imp}$ (Shephard's lemma). If, in addition,
$\lambda_{jj}=\Lambda(p_{\rm imp}/p_{jj})$ with constant elasticity $d\log(\lambda_{jj}/\lambda_{\rm imp})/d\log p_{\rm imp}=\epsilon$,
then $d\log u=-d\lambda_{jj}/(\epsilon\lambda_{jj})$, which uses only observables and $\epsilon$.
\end{proposition}
\begin{proof}
By Shephard's lemma, $d\log a=\sum_k\lambda_kd\log p_k$, and at fixed income $d\log u=-d\log a$ for homothetic $a$.
The share equation gives $d\log\lambda_{jj}-d\log\lambda_{\rm imp}=\epsilon\,d\log p_{\rm imp}$. Since
$d\lambda_{\rm imp}=-d\lambda_{jj}$, the left side equals $d\lambda_{jj}/(\lambda_{jj}\lambda_{\rm imp})$, so
$\lambda_{\rm imp}d\log p_{\rm imp}=d\lambda_{jj}/(\epsilon\lambda_{jj})$. Substitute into $d\log u=-\lambda_{\rm imp}d\log p_{\rm imp}$.
Cheaper imports lower $\lambda_{jj}$ and raise welfare, as in \cite{acr}.
\end{proof}

\section*{2. Non-identification for finite changes}
\begin{proposition}\label{prop:oe}
Fix baseline $(\lambda_{jj}^0,p^0)$ and local elasticity $\epsilon$. For a finite change $p_{\rm imp}^0\to p^1_{\rm imp}$ and
any value $\omega$ in an open interval around the CES welfare change, there is a homothetic expenditure function with the
same baseline shares and the same local elasticity at $p^0$, whose exact welfare change equals $\omega$.
\end{proposition}
\begin{proof}
Homothetic welfare is $-\int_{p^0}^{p^1}\lambda_{\rm imp}(p)\,d\log p_{\rm imp}$. Any smooth share function $\lambda_{\rm imp}(\cdot)$
with the given value and log-slope at $p^0$ is the import share of some homothetic preference, provided the implied
elasticity of substitution is nonnegative, i.e.\ $d\lambda_{\rm imp}/d\log p_{\rm imp}\le\lambda_{\rm imp}(1-\lambda_{\rm imp})$.
Then $\log e(p)=\int\lambda_{\rm imp}\,d\log p_{\rm imp}$ defines a concave unit-expenditure function. Start from the CES share,
which satisfies this strictly, and add a $C^1$-small bump supported away from $p^0$. Its integral can take any value in a small
interval, so it shifts the welfare integral by any amount in an open interval, while preserving the inequality, the baseline
share and the local elasticity.
\end{proof}
So structural gravity with a single elasticity is a local theory. Large counterfactuals, such as autarky gains, depend on the
global shape of import demand, which flows and a local elasticity do not reveal. Observing the share at several
price levels, i.e.\ historical episodes of large trade-cost changes, identifies the integral directly.

\section*{3. Nonhomotheticity}
\begin{proposition}\label{prop:nh}
Let households $h$ have nonhomothetic preferences, with import shares $\lambda^h_{\rm imp}$ depending on expenditure $m_h$.
Then to first order $\mathrm{EV}_h/m_h=-\lambda^h_{\rm imp}d\log p_{\rm imp}$, which is heterogeneous across $h$. Aggregate
flows identify only the expenditure-weighted mean $\bar\lambda_{\rm imp}$. Two economies with the same aggregate share and
elasticity but different covariances of $\lambda^h$ with $m_h$ have different distributions of EV.
\end{proposition}
\begin{proof}
Shephard's lemma household by household. The aggregate share is $\sum_hm_h\lambda^h/\sum m_h$.
\end{proof}
Household scanner or expenditure-survey data on import content by income identify $\lambda^h$, as in the source's
nonhomothetic measurement \cite{fk}.

\section*{4. Variable markups}
\begin{proposition}\label{prop:mk}
With markups $\mu_i=p_i/mc_i$ responding to trade costs, the first-order welfare change adds a term
$-\sum_i\lambda_i\,d\log\mu_i$ for domestically sold varieties, which is the pro-competitive effect, to the
cost-pass-through term. Flows and $\epsilon$ are invariant to how a given price change splits between $d\log mc$ and
$d\log\mu$. This split is therefore not identified from flows, but is identified from firm-level costs or markup estimates.
\end{proposition}
\begin{proof}
Prices enter welfare as above. Flows depend on prices only, while profits redistributed lump-sum change income by
$\sum\pi_i$, whose derivative contains the markup term.
\end{proof}

\section*{5. Remaining}
A multi-country general-equilibrium version requires income changes through wages, tariff revenue and profits. Our
results apply sector by sector, with income effects added. A full characterisation of the identified set of EV within a
specified $\Theta$ (for example all homothetic, quasi-separable demand systems with a given local elasticity) requires bounds
on the curvature of the import-share function. Such bounds turn Proposition~\ref{prop:oe}'s open interval into an explicit
interval, which we leave open \cite{av,chaney,fk}.

\begin{thebibliography}{9}
\bibitem{av} J. E. Anderson, E. van Wincoop, Gravity with gravitas, \emph{AER} 93 (2003), 170--192.
\bibitem{chaney} T. Chaney, Distorted gravity, \emph{AER} 98 (2008), 1707--1721.
\bibitem{fk} P. D. Fajgelbaum, A. K. Khandelwal, Measuring the unequal gains from trade, \emph{QJE} 131 (2016), 1113--1180.
\bibitem{acr} C. Arkolakis, A. Costinot, A. Rodr\'iguez-Clare, New trade models, same old gains?, \emph{AER} 102 (2012), 94--130.
\end{thebibliography}
\end{document}
